% =====================================================================
%  Introduction générale
% =====================================================================
\chapter*{Introduction générale}
\addcontentsline{toc}{chapter}{Introduction générale}
\markboth{Introduction générale}{}

En une quinzaine d'années, l'Internet des objets (\emph{Internet of Things}, IoT)
est passé du statut de concept à celui d'infrastructure du quotidien :
thermostats, montres, compteurs électriques, machines industrielles ou capteurs
agricoles sont désormais connectés. Le cabinet IoT Analytics estimait à environ
18,8~milliards le nombre d'objets connectés actifs dans le monde fin
2024~\cite{iotanalytics2024}. Ces objets produisent un volume de données
considérable, que l'architecture « classique » de l'IoT envoie vers le cloud pour
y être stockées, analysées et transformées en décisions.

Ce modèle centralisé montre aujourd'hui ses limites. Transmettre en permanence
des mesures brutes consomme de la bande passante et de l'énergie, introduit une
latence qui dépend de la qualité du réseau, rend le système inopérant lorsque la
connexion est coupée et expose des données parfois sensibles (sons, images,
mesures physiologiques). Le paradigme de l'\emph{edge computing} propose de
rapprocher le traitement de la source des données~\cite{shi2016edge}.

L'étape suivante consiste à doter les objets eux-mêmes d'une capacité
d'analyse : c'est l'\textbf{AIoT}, contraction d'\emph{Artificial Intelligence}
et d'\emph{Internet of Things}~\cite{zhang2021aiot}. Lorsque cette intelligence
est logée dans les composants les plus modestes de la chaîne, des
microcontrôleurs de quelques euros fonctionnant sur pile, on parle de
\textbf{TinyML}~\cite{warden2019tinyml,banbury2020benchmarking}. Un capteur
de vibrations ne transmet plus des milliers d'échantillons par seconde : il
signale simplement « fonctionnement anormal détecté ».

\begin{encadre}[title={Problématique}]
Comment faire fonctionner des modèles d'intelligence artificielle, conçus
habituellement pour des serveurs dotés de gigaoctets de mémoire, sur des objets
connectés qui ne disposent que de quelques centaines de kilooctets, et quels
bénéfices et limites cette approche apporte-t-elle aux systèmes IoT ?
\end{encadre}

Pour y répondre, ce rapport est organisé en cinq parties. La
\textbf{partie~1} pose les bases : objet connecté, architecture d'un système IoT,
définition de l'AIoT et du TinyML, et comparaison entre une architecture centrée
sur le cloud et une architecture où l'inférence est réalisée localement. La
\textbf{partie~2} détaille les principes techniques du TinyML, en particulier
le pipeline de développement et les techniques de compression des modèles. La
\textbf{partie~3} présente le matériel et la chaîne d'outils logiciels. La
\textbf{partie~4} examine les bénéfices de l'approche et ses domaines
d'application. Enfin, la \textbf{partie~5} discute les limites, les défis et les
perspectives d'avenir.
